\section*{Hoofdstuk \ref{ch:g6:light-propagation} --- Rechtlijnige voortplanting van licht}

\begin{solution}{exo:g6:light-propagation:1}
In heldere \omterm{def:g2:air-around-us:air}{lucht} (of in elk helder, goed gemengd materiaal) plant
\omterm{def:g1:light-and-shadows:source}{licht} zich langs rechte lijnen voort. Rechtlijnig: langs rechte
lijnen; voortplanting: het reizen van \omterm{def:g1:light-and-shadows:source}{licht} van plaats naar plaats.
\end{solution}

\begin{solution}{exo:g6:light-propagation:2}
Divergent: stralen die uit één punt uitwaaieren (een kaars, een bloot
\omterm{def:g3:first-electric-circuit:bulb}{lampje}). Evenwijdig: stralen naast elkaar (zonlicht --- de Zon staat
zo ver dat haar stralen ons in het gelid bereiken). Convergent:
stralen die naar een punt toe sluiten (\omterm{def:g1:light-and-shadows:source}{licht} gebundeld door een
vergrootinstrument).
\end{solution}

\begin{solution}{exo:g6:light-propagation:3}
Een \omterm{def:g6:light-propagation:ray}{lichtstraal} speelt de liniaal: oog, rand en verre uiteinde liggen
precies dan op één lijn wanneer \omterm{def:g1:light-and-shadows:source}{licht} van het verre uiteinde het oog
bereikt terwijl het langs de hele rand scheert --- de
\omterm{prop:g6:light-propagation:law}{rechtlijnige voortplanting} is de eed.
\end{solution}

\begin{solution}{exo:g6:light-propagation:4}
De landingspunten van de stralen die net langs de randen van de
blokkeerder scheren. Een oog in de schaduwkegel kan de lamp niet zien
--- geen enkele rechte straal vanaf de lamp bereikt het.
\end{solution}

\begin{solution}{exo:g6:light-propagation:5}
Halverwege: de \omterm{def:g1:light-and-shadows:shadow}{schaduw} is twee keer de bal --- \qty{12}{cm}. Dicht bij
het scherm: de \omterm{def:g1:light-and-shadows:shadow}{schaduw} krimpt naar de ware \qty{6}{cm} van de bal.
\end{solution}

\begin{solution}{exo:g6:light-propagation:6}
Onderaan het scherm: een rechte lijn vanuit een hoog punt door de lage
poort moet omlaag doorgaan.
\end{solution}

\begin{solution}{exo:g6:light-propagation:7}
Ondersteboven, en links-rechts verwisseld: op het scherm zwaait de
hand laag in plaats van hoog, en aan de andere kant. Beide omkeringen
komen van de draadjes die elkaar bij de poort kruisen.
\end{solution}

\begin{solution}{exo:g6:light-propagation:8}
Sterrenkundigen: zonsverduisteringen veilig op het scherm bekijken in
plaats van aan de hemel. Schilders: het geprojecteerde tafereel
overtrekken om het perspectief te beheersen. (Moderne verwanten: elke
camera, en het oog zelf.)
\end{solution}

\begin{solution}{exo:g6:light-propagation:9}
De \omterm{def:g1:light-and-shadows:shadow}{schaduw} krimpt naar de ware grootte van de munt: met de lamp
verder weg verlaten de scherende stralen haar minder steil --- een
smallere waaier --- en landen ze dichter bij elkaar. De constructie
laat het zien: vlakkere lijnen, strakkere spreiding.
\end{solution}

\begin{solution}{exo:g6:light-propagation:10}
De Zon is compact en ver: één strakke waaier bijna evenwijdige
stralen, één scherpe grens --- scherpe \omterm{def:g1:light-and-shadows:shadow}{schaduwen}. Een grijs
wolkendek is een lamp zo groot als de hele hemel: \omterm{def:g1:light-and-shadows:source}{licht} komt uit
alle richtingen tegelijk aan, elk punt van de grond wordt door een
deel ervan bereikt, en \omterm{def:g1:light-and-shadows:shadow}{schaduwen} spoelen weg tot vrijwel niets.
\end{solution}

\begin{solution}{exo:g6:light-propagation:11}
Helderder: een wijde poort laat van elk punt van het tafereel meer
draadjes toe. Onscherp: door een wijde poort bereiken de draadjes van
elk punt van het tafereel een hele \emph{vlek} scherm in plaats van
één plekje --- ontelbare overlappende vlekken smeren het \omterm{prop:g5:mirrors-reflection:image}{beeld} uit. De
scherpte van het gaatje \emph{is} zijn kleinheid.
\end{solution}

\begin{solution}{exo:g6:light-propagation:12}
Elk gaatje tussen de bladeren is een natuurlijke gaatjescamera; elke
lichte vlek op de grond is het \omterm{prop:g5:mirrors-reflection:image}{beeld} dat dat gaatje van de Zon maakt
--- rond op een gewone dag omdat de Zon rond is. Zodra de verduistering
de Zon in een sikkel verandert, projecteert elk bladgaatje getrouw de
sikkel: honderden kleine camera's, één hemels onderwerp.
\end{solution}

\begin{solution}{pb:g6:light-propagation:1}
\textbf{1.} Tegen het scherm hebben de scherende stralen geen afstand
meer om uit te waaieren: de \omterm{def:g1:light-and-shadows:shadow}{schaduw} sluit om de omtrek van de pop zelf.
\textbf{2.} Halverwege verdubbelt hem: \qty{40}{cm}.
\textbf{3.} Vier keer $20$ is precies \qty{80}{cm} --- het klopt; de
pop staat op een kwart van de afstand lamp--scherm.
\textbf{4.} Een brede gloeiende bol is vele lampen tegelijk: elk deel
werpt een eigen, iets verschoven \omterm{def:g1:light-and-shadows:shadow}{schaduw}, en de randen van de draak
smeren zacht en grijs uit. Scherp theater heeft een kleine, blote,
puntvormige bron nodig.
\textbf{5.} Naar de lamp toe: de waaier van scherende stralen wordt
wijder en de wolf groeit; de arm blijft laag, buiten de waaier.
Naarmate de pop de kleine maar echte lamp nadert, groeit de
randverzachting een beetje --- van dichtbij telt de breedte van de
lamp zwaarder.
\textbf{6.} De rechte lijnen van de lamp door de twee poppen moeten
naar aparte gebieden van het scherm lopen --- de poppen staan
zijdelings ver genoeg uit elkaar dat geen van beide in de
schaduwkegel van de ander staat.
\textbf{7.} Die van de hand: dichter bij de lamp, op een derde van de
weg, waaieren haar scherende stralen drie keer zo wijd uit, terwijl
die van de draak halverwege maar verdubbelen --- vlak bij de lamp
spelen zelfs kleine handen voor reus.
\textbf{8.} Beweeg de schurk plat tegen het scherm en hij krimpt tot
zijn gewone zelf (of helemaal uit de \omterm{def:g6:light-propagation:ray}{bundel}); of schuif een tweede
blokkeerder dicht bij de lamp --- zijn reuzenschaduw slokt die van de
schurk op, en de schurk ``verdwijnt'' in een andere duisternis.
\textbf{9.} Laag op het laken: vanaf de hoge lamp gaat het rechte
draadje door het gaatje omlaag verder --- het gaatje keert om.
\textbf{10.} Van rechts naar links: de poort verwisselt links en
rechts, net als boven en onder.
\textbf{11.} Helderder maar hopeloos onscherp: elk punt van de straat
schildert nu een muntgrote vlek, de vlekken overlappen, en het \omterm{prop:g5:mirrors-reflection:image}{beeld}
lost op --- de finale sterft met de scherpte.
\textbf{12.} ``Het hele programma van vanavond --- elke
\omterm{def:g1:light-and-shadows:shadow}{schaduw}, elke reus en één omgekeerde straat --- is uitgevoerd door
één enkele wet: \omterm{def:g1:light-and-shadows:source}{licht} plant zich langs rechte lijnen voort.''
\end{solution}
